\begin{statement}{C.2.UN.59}{Proved}{All multiplicities and all diagonal Hopf and ruled vertices}
Specify a simple-NC cycle of smooth embedded marked surfaces, with
edges between distinct components. Each vertex is a diagonal primary
or secondary Hopf surface with its two axes, or
$\PP(\OO_E\oplus L)$ with $E$ elliptic and
$L\in\operatorname{Pic}^0(E)$, marked by its two sections.
Specify their elliptic gluing maps, positive multiplicities $m_i$,
and the actual compatible weighted normal lines, sections and
boundary maps of (C.WH.1)--(C.WH.2).

Put $N=\operatorname{lcm}_i m_i$. Form the actual normalized
$N$th-root component covers and the compressed necks $uv=s^{h_i}$.
The root graph is connected and its rank-three integral
complement-lattice product is the identity precisely when the
specified weighted normal plumbing has connected two-torus fibres.
Every actual smooth-total torus degeneration with the specified
marked support yields these data and passes this test.

Every weighted support passing the test occurs over $\PP^1$ in a
smooth compact complex threefold. The construction retains all
marked components, elliptic gluing maps and multiplicities. Its
only other singular fibre is the complex conjugate weighted
support. All reduced fibre supports have Euler number zero.
Normal-product constants are chosen in this compact realization.
There is no algebraic-dimension restriction on the Hopf vertices,
and an all-ruled cycle is included.

The complete component, axis and inertia comparisons in
(C.WH.0), and the ruled-cover identities (C.WM.2)--(C.WM.7),
hold in this scope. In common root-complement coordinates the
finite deck linear part fixes the primitive total normal period.
Its other eigenvalue is a root of unity of order $1,2,3,4$ or $6$.
The compact construction retains this linear part and every
affine deck constant.
\end{statement}

\begin{proof}
\par\medskip\noindent\textbf{Actual root maps and the local criterion.}
The weighted line on a diagonal Hopf vertex has a flat tensor
power, since both axis divisors have the explicit constant
automorphy factors of UN40 and UN43. Dividing the transferred
connection by its tensor exponent gives the actual flat line
used in UN52. Its root cover on the axis complement is the
finite-index logarithmic lattice cover constructed there.
Coordinate powers extend it across both axes, and the actual
contracting generator and free finite angular characters identify
its finite normalization with a smooth diagonal Hopf surface.
This comparison does not require nonresonant multipliers.

On a ruled vertex the proof (C.WM.4)--(C.WM.7) applies without
any Hopf assumption. It gives the actual finite normalization
$\PP(\OO_{E'}\oplus L')\to\PP(\OO_E\oplus L)$, with
an \'{e}tale base map $\pi$, fibre power $q$, and
$(L')^q\simeq\pi^*L$. In particular
$m\mid a+b$ and $q=m/\gcd(m,a)=m/\gcd(m,b)$.
The weighted normal line is pulled back on each entire section
chart by its transferred flat connection; its nonzero global
degree on a ruling fibre is retained as in UN56.

At an original double curve put
$a=gA$, $b=gB$, $\gcd(A,B)=1$ and
$h=N/(gAB)$. The finite normalization is still exactly
\[
             uv=s^h,\qquad x=u^B,\quad y=v^A,
             \qquad t=s^N=x^a y^b.
\tag{C.EW.1}
\]
The actual root component and axis covers, their connected
degrees, and all stabilizers therefore have precisely the
maps and orders proved in UN52 and UN56. Each strict component
has two connected marked elliptic curves, so the compressed
root graph is a union of cycles. Resolving a neck inserts the
specified degree-zero ruled bridges and changes no complement
lattice product.

Connectedness of this graph is equivalent to connectedness of
the root family, by the proper Stein factor and its reduced
central fibre. On a connected graph the actual core-and-neck
calculation gives $\pi_1(F_s)=\Lambda\rtimes_M\Z$.
Thus a torus forces $M=I$. Conversely $M=I$ glues the closed
logarithmic coframe; its complete commuting fields identify
the nearby compact fibre with $\C^2$ modulo its rank-four
period lattice. For an arbitrary actual smoothing the
canonical divisor sections and their small smooth unit roots
give the same weighted tubular and collar maps, exactly as
in UN52. This proves the asserted local criterion for the
actual maps, as well as for the explicit plumbing.

\par\medskip\noindent\textbf{The exact finite linear and affine data.}
Assume that the test passes and identify all open root
complements with $G=\C^2/\Lambda$ by the actual linear
edge maps. Let $n_j$ be the outgoing normal meridian at
vertex $j$, and let $h_j$ be its outgoing compressed
neck exponent. The incoming meridian is $-n_{j-1}$.
The root generator lies over the identity on each original
marked normalization and preserves its two marked branches.
It consequently induces an orientation-preserving rotation
$\sigma$ of the connected root cycle. The common number
$r$ of strict covers and every inertia order are therefore
exactly those of (C.WH.0).

The generator is affine on the common complement. For clarity,
an isomorphism of diagonal Hopf surfaces preserving both axes
lifts to a biholomorphism of $\C^2\setminus0$, extends with
its inverse across zero, and has coordinates $zU(z,w)$ and
$wV(z,w)$, with $U,V$ units. Conjugating a contracting deck
generator and evaluating at zero makes the corresponding
multiplier of each unit equal to one. Iteration of the
contraction then makes each unit constant. Thus this lift
is diagonal even for resonant multipliers. For a ruled
vertex the actual base isomorphism and marked line-bundle
isomorphism have affine logarithmic lifts by the constant
factors of automorphy, as in UN51. Compatibility with the
actual neck maps makes the linear part a single map $R$.
Write the generator and the edge maps as
\[
 \begin{aligned}
 g(j,v,s)&=(\sigma j,Rv+\tau_j,\zeta_Ns),\\
 e_j(v,s)&=v+b_j+l_j\log s,\\
 l_j&=-\frac{h_j n_j}{2\pi i}.
 \end{aligned}
\tag{C.EW.2}
\]
Here $R\Lambda=\Lambda$, $R^N=I$,
$Rn_j=n_{\sigma j}$ and $h_{\sigma j}=h_j$.
Actual composition of these maps, with
$\log\zeta_N=2\pi i/N$, gives lattice elements
$\lambda_j,\mu_j\in\Lambda$ satisfying
\[
 \begin{aligned}
 \tau_{j+1}+Rb_j-\tau_j-b_{\sigma j}
             +\frac{h_j n_{\sigma j}}N&=\lambda_j,\\
 \sum_{a=0}^{N-1}R^{N-1-a}\tau_{\sigma^a j}&=\mu_j.
 \end{aligned}
\tag{C.EW.3}
\]
These are equalities in the actual logarithmic vector space.

Put
\[
             b=\sum_jb_j,\qquad
             u=-\sum_jh_jn_j=d p\in\Lambda,
\tag{C.EW.4}
\]
where $p$ is primitive and $d>0$. The contraction-height
calculation (C.AD.3) applies at every diagonal Hopf vertex;
at a ruled vertex the consecutive outgoing normals agree.
Thus for $W=\operatorname{span}_{\R}\Lambda$ and
$\ker\ell=W$, all numbers $\ell(i n_j)$ have one
nonzero sign. In the all-ruled case the same assertion
is the actual elliptic quotient calculation of UN55.
Consequently $\ell(iu)\ne0$. Rotation and the displayed
normal maps give $Ru=u$, hence $Rp=p$.

The complex line $W\cap iW$ is $R$-invariant and is
complementary to $\C p$. Choose the complex coordinates
given by these two lines, with $p=(0,1)$. Then
\[
 \begin{gathered}
 W=\C\times\R,\qquad R=\operatorname{diag}(\rho,1),\\
 \Lambda=\langle w_1=(a_1,c_1),\ w_2=(a_2,c_2),\ p\rangle_{\Z},\\
 c_1,c_2\in\R,\qquad
 C=\operatorname{Im}(\overline{a_1}a_2)>0.
 \end{gathered}
\tag{C.EW.5}
\]
The projection of $\Lambda$ to the first coordinate is
the actual elliptic lattice. Since $R$ preserves it,
$\rho$ has order $1,2,3,4$ or $6$: its integral real
two-dimensional matrix has determinant one, trace
$2\operatorname{Re}\rho\in\Z$, and finite order.
Also $\rho^N=1$. The stabilizer of a root vertex acts
as the deck translations of its finite open-group cover,
so in fact $R^r=I$.

Summing (C.EW.3) retains the full affine descent relation
\[
                         (R-I)b-u/N\in\Lambda.
\tag{C.EW.6}
\]
When $R=I$ this recovers $u/N\in\Lambda$ from UN52.
Equation (C.EW.6), with its possibly nontrivial $R$, is
the relation used below.

Adding $c n_j$ to $b_j$ on a whole $\sigma$-orbit,
with the same $c\in\C$, changes only the normal-product
constants. It preserves both equations (C.EW.3), since
$Rn_j=n_{\sigma j}$. The second coordinates of all
$n_j$ are real and have the same nonzero sign, by the
height calculation. The sum on any such orbit therefore
has nonzero second coordinate. Choose $c$ on one orbit
so that
\[
                              b=(\sigma_0,i).
\tag{C.EW.7}
\]
Every original elliptic gluing map and every finite
action map on the central support is retained.
The orbit-normal sum is also $R$-fixed; when
$\rho\ne1$ its first coordinate is zero. Thus the
entire common horizontal affine constant is retained
in that case as well.

\par\medskip\noindent\textbf{A finite-equivariant period family on the punctured sphere.}
Choose $k\in\{1,\ldots,N\}$ with
$\rho=\zeta_N^k$, using $k=N$ for $\rho=1$.
Let
\[
 J=\begin{pmatrix}0&1\\-1&0\end{pmatrix},\qquad
 V_s(w)=w+s^kJ\overline w.
\tag{C.EW.8}
\]
For each fixed vector $w$ this is holomorphic in $s$.
The exact relation $RJ\overline R^{-1}=\rho J$
gives
\[
                           R V_s=V_{\zeta_Ns}R.
\tag{C.EW.9}
\]
Transport all integral vectors by $V_s$, and define
\[
 \begin{gathered}
 \Lambda_s=V_s\Lambda,\qquad
 B_s=V_s b+\frac{d\log s}{2\pi i}V_s p,\\
                         T_s=\C^2/(\Lambda_s+\Z B_s).
 \end{gathered}
\tag{C.EW.10}
\]
A logarithm change adds $dV_s p$, an actual integral
period. The determinant calculation of UN55 with
$z=s^k$ and logarithmic coefficient $d\log s/(2\pi i)$
gives, for $r_s=|s|>0$,
\[
 C(1+r_s^{2k})
 \left(1+r_s^{2k}
       +\frac d{2\pi}(r_s^{2k}-1)\log r_s\right)>0.
\tag{C.EW.11}
\]
Thus these are actual compact two-tori for every
$s\in\C^*$, and form a smooth proper holomorphic
family there. Properness follows from the uniformly
bounded period matrices and inverses on compact
base charts, exactly as in UN55.

The transported finite action and edge data are explicitly
\[
 \begin{aligned}
 b_j(s)&=V_s b_j,& n_j(s)&=V_s n_j,\\
 l_j(s)&=-h_jn_j(s)/(2\pi i),&
 \tau_j(s)&=V_{\zeta_Ns}\tau_j.
 \end{aligned}
\tag{C.EW.12}
\]
Notice that the translation is expressed in the
\emph{target} frame. Applying (C.EW.9) to the first
equation of (C.EW.3) gives exactly
\[
 \begin{aligned}
 &\tau_{j+1}(s)+Rb_j(s)-\tau_j(s)-b_{\sigma j}(\zeta_Ns)\\
 &\qquad-l_{\sigma j}(\zeta_Ns)\,2\pi i/N
       =V_{\zeta_Ns}\lambda_j.
 \end{aligned}
\tag{C.EW.13}
\]
The logarithm correction transports because
$h_jn_{\sigma j}/N$ is a real rational multiple
of an integral vector. For the order equation,
\[
 R^{N-1-a}V_{\zeta_N^{a+1}s}
                         =V_sR^{N-1-a},
\]
so the $N$th iterate translates by $V_s\mu_j$.
This proves the full finite action, including all
overlap and order equalities. Likewise
\[
                    R B_s-B_{\zeta_Ns}
                   =V_{\zeta_Ns}\bigl((R-I)b-u/N\bigr)
                   \in\Lambda_{\zeta_Ns}.
\tag{C.EW.14}
\]

Here is an explicit global comparison with (C.EW.10).
On a logarithm chart put
$\theta_j(s)=\sum_{a<j}(b_a(s)+l_a(s)\log s)$.
The common coordinate of the $j$th component chart is
$v_j-\theta_j(s)$. The generator on the period quotient
therefore has the affine translation
\[
 R\theta_j(s)+\tau_j(s)-\theta_{\sigma j}(\zeta_Ns).
\tag{C.EW.15}
\]
Changing $j$ changes this by an integral period, by
(C.EW.13); passing the last edge adds exactly $B_s$.
Changing the logarithm changes each partial sum by
an integral vector. Thus (C.EW.15) defines a single
holomorphic finite action on the entire punctured
period family. Its order is the order already proved
in (C.EW.13), and its linear and period maps are
(C.EW.9) and (C.EW.14).

\par\medskip\noindent\textbf{The actual finite filling at zero.}
Near zero construct the moving strict components from
their transported lattices. At a Hopf vertex use
\[
 P_j(s)=\left(\frac{V_s n_j}{2\pi i},
                        -\frac{V_s n_{j-1}}{2\pi i}\right).
\tag{C.EW.16}
\]
It is invertible near zero, retains the distinguished
angular periods and their finite angular enlargement,
and transports the original contracting logarithms
holomorphically. Their real parts remain negative.
It therefore defines the actual marked Hopf family.
At a ruled vertex use its primitive normal and an
independent projected elliptic lattice lift, as in
UN56 and UN57. This gives the actual moving
degree-zero line and its projective completion.

Use the edge constants (C.EW.12) and keep the literal
compressed hypersurfaces $uv=s^{h_j}$. They give the
proper moving plumbing, whose exact periods are
(C.EW.10). The component logarithmic maps and their
closed-coframe inverses give a holomorphic
biholomorphism of its entire punctured family with
that period family. Formula (C.EW.15) makes this
biholomorphism equivariant.

The finite maps extend across all marked curves.
At a Hopf vertex
$RP_j(s)=P_{\sigma j}(\zeta_Ns)$, so the logarithmic
action in these component coordinates is a translation;
its exponential is the actual diagonal map on the
marked Hopf family. At a ruled vertex the same
normal-period map and its elliptic base map give the
actual line-bundle map and homogeneous projective
extension from UN56. The overlap equalities
(C.EW.13) extend them across the compressed necks.

The smoothness of the finite quotient is the same
local ring calculation as (C.WH.12), which uses
the actual inertia characters rather than $R=I$.
At (C.EW.1) the pointwise axis stabilizer has order
$hAB$. Its order-$h$ subgroup fixes $u,v$ and rotates
$s$; the remaining independent characters have
orders $B,A$. These finite characters retain their
specified central values throughout the moving
families. We give the actual character coordinates.
The stabilizer preserves each adjacent component,
so it sends each of $u,v$ to itself times a unit.
Reynolds-average these functions to their specified
central eigencharacters. The results are $u'=uU$
and $v'=vV$, with $U,V$ units. Average the coordinate
along the pointwise fixed axis to an invariant
coordinate; its derivative there is unchanged.
The character identity
$\chi_u\chi_v=\chi_s^h$ and the equation
$u'v'=s^hUV$ show that $A=UV$ is invariant.
Replace $u'$ by $u'/A$. The product is then exactly
$s^h$, both eigencharacters are retained, and the
actual base coordinate $s$ has not changed. The
coordinate change has invertible differential at
the axis and is an analytic local isomorphism.
In these exact coordinates the ring is
\[
 \bigl(\C\{u,v,s\}/(uv-s^h)\bigr)^{J_{\rm axis}}
                         =\C\{u^B,v^A\}
                         =\C\{x,y\},
 \qquad t=x^a y^b.
\tag{C.EW.17}
\]
The residual axis cover is \'{e}tale. At a general
component point the inertia is the normal reflection
of order $N/m_i$. Thus the quotient is smooth at
every point and has exactly the original full
multiplicities. At zero the covers, central maps and
finite action are the original ones; the quotient
support is therefore the specified marked surface.

\par\medskip\noindent\textbf{The conjugate finite filling at infinity.}
Put $w=s^{-1}$ and divide the vector frame by $s^k$.
For a fixed vector $c$ write
\[
 W_w(c)=J\overline c+w^k c.
\tag{C.EW.18}
\]
The periods in this frame are $W_w\Lambda$ and
\[
                    W_wb-\frac{d\log w}{2\pi i}W_wp.
\tag{C.EW.19}
\]
Choose outgoing normals $-W_wn_j$, the same neck
exponents $h_j$, and affine edge constants $W_wb_j$.
Their logarithmic coefficients are
$h_jW_wn_j/(2\pi i)$, so their complete circuit is
exactly (C.EW.19).

The actual transformed deck generator is
\[
 \begin{aligned}
 (j,v,w)&\longmapsto
     (\sigma j,R_\infty v+\tau_{j,\infty}(w),\zeta_N^{-1}w),\\
 R_\infty&=\rho^{-1}R=J\overline R J^{-1},\\
 \tau_{j,\infty}(w)
     &=J\overline{\tau_j}+\rho^{-1}w^k\tau_j
       =W_{\zeta_N^{-1}w}\tau_j.
 \end{aligned}
\tag{C.EW.20}
\]
Indeed these formulas follow by multiplying the image
vector $Rv+\tau_j(s)$ by $(\zeta_Ns)^{-k}$.
They retain all finite equations, since they are the
actual frame change of the action already constructed.

At $w=0$ the Hopf frame is exactly
$J\overline{P_j(0)}$, as in (C.AD.11).
The angular normals have the displayed minus sign
because conjugation changes $2\pi i$ to $-2\pi i$.
Thus the contracting logarithms and every finite
angular character are their actual complex
conjugates. The separate ruled frame gives the
conjugate line bundle and marked sections by the
same logarithmic calculation. Each affine edge
constant, in its component coordinates, is the
conjugate of the original one. Equation (C.EW.20)
is precisely the conjugate central action under
the anti-linear comparison $v\mapsto J\overline v$,
with the conjugate base rotation.

The resulting proper compressed filling at infinity
is therefore the complex conjugate root plumbing
with its actual finite action. Its punctured-family
map has exactly (C.EW.19) and (C.EW.20).
Its quotient rings are the conjugates of
(C.EW.17), so the quotient is smooth and its full
central fibre has the conjugate support with the
same multiplicities.

Finally glue the two finite-equivariant proper
fillings to the entire period family (C.EW.10),
using their explicit logarithmic biholomorphisms
and the frame change (C.EW.18). The closed-base
refinement argument of UN55 proves this normal
root total space Hausdorff and proper over
$\PP^1_s$. Take the actual finite quotient by
the constructed action. The base quotient is
$\PP^1_t$ with $t=s^N$. Its total space is smooth
by (C.EW.17) at the two special values and by
freeness of the base action elsewhere. It is
compact and proper. A general quotient fibre is
identified with any one of its $N$ root fibres,
so is the specified connected two-torus.
Every component and every double curve of either
special support has Euler number zero. Additivity
therefore proves the final Euler assertion.

\par\medskip\noindent\textbf{Explicit nontrivial finite linear actions.}
For $N=2,3,4,6$, choose an elliptic curve $E$ with
an origin-preserving automorphism $\rho$ of exact
order $N$. Start with a root cycle of $2N$ copies
of $E\times\PP^1$, glued by the identity on their
elliptic sections. In the coordinates above take
\[
 \begin{gathered}
 \Lambda=\Lambda_E\oplus\Z p,\quad
 n_j=-p,\quad h_j=1,\quad
 b_j=ip/(2N),\quad \tau_j=j p/N,\\
 \sigma(j)=j+2\pmod{2N},\qquad
 R=\operatorname{diag}(\rho,1).
 \end{gathered}
\tag{C.EW.21}
\]
The first equation (C.EW.3) is zero except at the
cyclic endpoint, where it is $-2p$.
The $N$th translation sum is $(N-1)p$ on even
vertices and $Np$ on odd vertices. These are
actual integral periods, so this is the specified
finite action. It is free on the central graph.
The quotient has two components $E\times\PP^1$,
both with full multiplicity $N$. One elliptic
gluing is the identity; the other is $\rho^{-1}$,
as follows by identifying vertex $2$ with vertex
$0$ using the inverse generator. The root cycle
has identity product, while this latter gluing
retains the nontrivial finite linear part.
The compact construction above realizes these
exact weighted supports and their conjugates.

The equivariance, target-frame affine equations,
infinity identities, positive determinant and
all four finite-action examples are checked
exactly in \path{code/D_all_weighted_cycles.py}
and its saved output.
\end{proof}


% C.2.UN.17 from part_C2_canonical_components.tex

